% =============================================================================
% Chapter 3: Limits of Classical Utility
% =============================================================================
%
% Version: 1.0 (January 2026)
%
% Purpose: Why extensions needed
%
% Primary Appendix: See appendix references below
% Secondary Appendices: G (Glossary), F (Examples)
%
% Prerequisites: None
% Leads to: Chapter 4
%
% Chapter Type: B (Foundation)
% =============================================================================


%%%%%%%%%%%%%%%%%%%%%%%%%%%%%%%%%%%%%%%%%%%%%%%%%%%%%%%%%%%%%%%%%%%%%%%%%%%%%%%
% QUICK REFERENCE BOX
%%%%%%%%%%%%%%%%%%%%%%%%%%%%%%%%%%%%%%%%%%%%%%%%%%%%%%%%%%%%%%%%%%%%%%%%%%%%%%%
\begin{tcolorbox}[
    colback=gray!5!white,
    colframe=gray!75!black,
    title={\textbf{Begriffe in diesem Kapitel}},
    fonttitle=\bfseries
]
\small
\begin{tabular}{@{}ll@{}}
\textbf{Key Concept} & \textbf{Definition} \\
\midrule
Why extensions needed & See Section \ref{sec:ch3-intro} \\
\end{tabular}
\end{tcolorbox}



%%%%%%%%%%%%%%%%%%%%%%%%%%%%%%%%%%%%%%%%%%%%%%%%%%%%%%%%%%%%%%%%%%%%%%%%%%%%%%%
% INTUITION BOX
%%%%%%%%%%%%%%%%%%%%%%%%%%%%%%%%%%%%%%%%%%%%%%%%%%%%%%%%%%%%%%%%%%%%%%%%%%%%%%%
\begin{tcolorbox}[
    colback=blue!5!white,
    colframe=blue!75!black,
    title={\textbf{Intuition: A Motivating Example}},
    fonttitle=\bfseries
]
Consider \textbf{Anna}, a professional facing a decision that involves why extensions needed.

She initially evaluates the choice using standard criteria, but something feels incomplete.
\textbf{Thomas}, her colleague, suggests considering additional dimensions beyond the obvious.

This chapter explores what Anna discovers when she applies the EBF framework---how
why extensions needed interact with broader welfare considerations.

\medskip
\textit{By the end of this chapter, you will understand why Anna's initial analysis
was incomplete and how the EBF approach provides a more comprehensive view.}
\end{tcolorbox}



%%%%%%%%%%%%%%%%%%%%%%%%%%%%%%%%%%%%%%%%%%%%%%%%%%%%%%%%%%%%%%%%%%%%%%%%%%%%%%%
% CENTRAL QUESTION BOX
%%%%%%%%%%%%%%%%%%%%%%%%%%%%%%%%%%%%%%%%%%%%%%%%%%%%%%%%%%%%%%%%%%%%%%%%%%%%%%%
\begin{tcolorbox}[
    colback=red!5!white,
    colframe=red!75!black,
    title={\textbf{Central Question}},
    fonttitle=\bfseries
]
\Large
\centering
\textit{How does why extensions needed contribute to understanding human welfare and decision-making?}
\end{tcolorbox}



%%%%%%%%%%%%%%%%%%%%%%%%%%%%%%%%%%%%%%%%%%%%%%%%%%%%%%%%%%%%%%%%%%%%%%%%%%%%%%%
% CHAPTER OVERVIEW
%%%%%%%%%%%%%%%%%%%%%%%%%%%%%%%%%%%%%%%%%%%%%%%%%%%%%%%%%%%%%%%%%%%%%%%%%%%%%%%
\subsection*{Chapter Overview}

This chapter is organized as follows:

\begin{itemize}
\item \textbf{Section \ref{sec:ch3-intro}}: Introduction and motivation
\item \textbf{Section \ref{sec:ch3-theory}}: Theoretical foundations
\item \textbf{Section \ref{sec:ch3-applications}}: Applications and examples
\item \textbf{Section \ref{sec:ch3-summary}}: Summary and conclusions
\end{itemize}




%%%%%%%%%%%%%%%%%%%%%%%%%%%%%%%%%%%%%%%%%%%%%%%%%%%%%%%%%%%%%%%%%%%%%%%%%%%%%%%
% CHAPTER SCOPE BOX
%%%%%%%%%%%%%%%%%%%%%%%%%%%%%%%%%%%%%%%%%%%%%%%%%%%%%%%%%%%%%%%%%%%%%%%%%%%%%%%
\begin{tcolorbox}[
    colback=purple!5!white,
    colframe=purple!75!black,
    title={\textbf{Chapter Scope: Ziel / In-Scope / Out-of-Scope}},
    fonttitle=\bfseries
]

\textbf{Ziel (Objective):}
Develop the conceptual understanding of why extensions needed as a foundation for the EBF framework.

\vspace{0.3cm}
\textbf{In-Scope:}
\begin{itemize}[nosep]
    \item Conceptual introduction to why extensions needed
    \item Motivation and intuition
    \item Connection to subsequent chapters
    \item Key definitions and concepts
\end{itemize}

\vspace{0.3cm}
\textbf{Out-of-Scope (delegiert an Appendices):}
\begin{itemize}[nosep]
    \item Complete formal treatment $\rightarrow$ FORMAL appendices
    \item Empirical validation $\rightarrow$ METHOD appendices
    \item Domain applications $\rightarrow$ Application chapters
\end{itemize}

\end{tcolorbox}

%%%%%%%%%%%%%%%%%%%%%%%%%%%%%%%%%%%%%%%%%%%%%%%%%%%%%%%%%%%%%%%%%%%%%%%%%%%%%%%
% APPENDIX REFERENCES BOX
%%%%%%%%%%%%%%%%%%%%%%%%%%%%%%%%%%%%%%%%%%%%%%%%%%%%%%%%%%%%%%%%%%%%%%%%%%%%%%%
\begin{tcolorbox}[
    colback=yellow!10!white,
    colframe=orange!75!black,
    title={\textbf{Appendix References for This Chapter}},
    fonttitle=\bfseries
]
\small
\begin{tabular}{@{}lp{8cm}@{}}
\textbf{Appendix} & \textbf{Content} \\
\midrule
A (FORMAL-LIMITS) & Formal derivations of behavioral limiting cases \\
U (LIT-KAHNEMAN) & Kahneman-Tversky prospect theory foundations \\
G (REF-GLOSSARY) & Symbol definitions and terminology \\
\end{tabular}

\medskip
\textit{For formal proofs and complete derivations, see the referenced appendices.}
\end{tcolorbox}

\section{The Limits of Utility Maximization under Stable Preferences}
\label{sec:limitsutility}
%%%%%%%%%%%%%%%%%%%%%%%%%%%%%%%%%%%%%%%%%%%%%%%%%%%%%%%%%%%%%%%%%%%%%%%%%%%%%%%

The classical framework assumes that agents maximize well-defined utility functions
with stable preferences. This section examines the empirical and theoretical challenges
to this assumption, demonstrating that the stability assumption is not merely
a simplification but a substantive restriction that excludes central economic phenomena. Formal derivations showing how each behavioral phenomenon emerges as a limiting case appear in Appendix~A.

\subsection{The Stable Preferences Assumption}

\paragraph{What the Assumption Claims.}
Stable preferences means:
\begin{enumerate}
  \item Utility functions don't change over time: $U_i(x, t) = U_i(x)$
  \item Preferences are independent of others: $U_i(x_i, x_j) = U_i(x_i)$
  \item Context doesn't affect valuation: $U_i(x; \Psi) = U_i(x)$
\end{enumerate}

\paragraph{Why Economists Made This Assumption.}
The assumption is not arbitrary---it serves crucial theoretical functions:
\begin{itemize}
  \item \textbf{Tractability:} Fixed preferences enable optimization
  \item \textbf{Prediction:} Stable preferences allow demand estimation
  \item \textbf{Welfare:} Unchanging preferences define a welfare criterion
  \item \textbf{Identification:} Without stability, any behavior is ``rational''
\end{itemize}

\paragraph{Stigler-Becker Defense.}
\citet{stiglerbecker1977} argued that apparent preference changes can always
be reinterpreted as responses to changing constraints.
``De gustibus non est disputandum''---tastes don't change; only opportunities do.

\paragraph{Example: Addiction.}
Is addiction a changing preference or a stable preference with habit formation?

Stigler-Becker interpretation:
\begin{equation}
  U(c_t, S_t) \quad \text{where } S_t = \sum_{\tau < t} c_\tau \text{ is consumption stock}
\end{equation}

The utility function is stable; past consumption enters as a state variable.

\paragraph{The Problem.}
This reinterpretation saves stability but at a cost:
it makes the theory unfalsifiable.
Any behavior pattern can be rationalized with a sufficiently complex constraint structure.

\subsection{Six Domains Where Stable Preferences Fail}

\subsubsection{Domain 1: Social Preferences}

\paragraph{The Evidence.}
Thousands of experiments show that people care about others' outcomes:
\begin{itemize}
  \item \textbf{Ultimatum games:} Responders reject unfair offers, sacrificing money \citep{guth1982}
  \item \textbf{Dictator games:} Dictators share even when anonymous
  \item \textbf{Public goods games:} Many contribute despite dominant strategy to free-ride \citep{andreoni1988}
  \item \textbf{Trust games:} Investors trust and trustees reciprocate \citep{berg1995}
\end{itemize}

\paragraph{Example: The Ultimatum Game.}
Proposer divides \$10. Responder accepts or rejects. If reject, both get zero.

\textbf{Classical prediction:} Proposer offers \$0.01; responder accepts.

\textbf{Actual behavior:} Modal offer is \$4-5; offers below \$2 often rejected.

\paragraph{Implication.}
Utility is not $U_i(x_i)$ but $U_i(x_i, x_j)$.
The cross-derivative is non-zero:
\begin{equation}
  \frac{\partial^2 U_i}{\partial x_i \partial x_j} \neq 0
\end{equation}

This is precisely the complementarity that classical theory assumes away.

\paragraph{Models of Social Preferences.}
Several models formalize other-regarding preferences:
\begin{itemize}
  \item \textbf{Inequality aversion} \citep{fehrschmidt1999}: 
        $U_i = \pi_i - \alpha \max\{\pi_j - \pi_i, 0\} - \beta \max\{\pi_i - \pi_j, 0\}$
  \item \textbf{ERC model} \citep{boltonockenfels2000}: Utility depends on own payoff and relative share
  \item \textbf{Social welfare} \citep{charnessrabin2002}: Combines efficiency and maximin concerns
  \item \textbf{Reciprocity} \citep{rabin1993}: Kindness begets kindness; unkindness begets punishment
\end{itemize}

\paragraph{The State of the Art.}
The comprehensive survey by \citet{fehrcharness2024} synthesizes three decades of
research on social preferences, documenting:
\begin{enumerate}
  \item \textbf{Distributional preferences:} Inequality aversion and altruism
  \item \textbf{Belief-dependent preferences:} Reciprocity and guilt aversion \citep{charnessdufwenberg2006}
  \item \textbf{Image concerns:} Self-image and social image \citep{benaboutirole2006}
\end{enumerate}

\paragraph{Distribution of Types.}
A key finding: the majority of individuals have social preferences;
purely selfish individuals are a minority:
\begin{center}
\begin{tabular}{lcc}
\toprule
\textbf{Type} & \textbf{Students} & \textbf{General Population} \\
\midrule
Selfish & 40--60\% & 20--30\% \\
Inequality averse & 10--15\% & 15--20\% \\
Altruistic & 25--45\% & 45--55\% \\
\bottomrule
\end{tabular}
\end{center}

\paragraph{Merit and Desert.}
Social preferences are context-dependent. Tolerance for inequality depends on
its source \citep{cappelen2007,almas2020}:
\begin{itemize}
  \item Inequality from effort/merit: More accepted
  \item Inequality from luck: Less accepted, more redistribution demanded
  \item Cultural variation: Americans more meritocratic, Scandinavians more egalitarian
\end{itemize}

\paragraph{Context Dependence.}
The magnitude of social preferences varies with context:
\begin{itemize}
  \item Higher in face-to-face interactions
  \item Higher with in-group members
  \item Lower in competitive frames
  \item Lower when ``deserving'' vs. ``undeserving'' recipients
\end{itemize}

Thus: $C_{ij} = C_{ij}(\Psi)$---complementarity depends on context.

\subsubsection{Domain 2: Identity and Self-Image}

\paragraph{The Evidence.}
People sacrifice material payoffs to maintain consistent self-image:
\begin{itemize}
  \item \textbf{Professional identity:} Doctors resist cost-cutting that conflicts with medical ethics
  \item \textbf{Group identity:} Fans support losing teams; patriots sacrifice for country
  \item \textbf{Moral identity:} People avoid situations that would reveal their selfishness
\end{itemize}

\paragraph{Example: The Moral Wiggle Room Experiment.}
\citet{danaberton2007} offered subjects a choice:
\begin{itemize}
  \item Option A: Get \$6, other gets \$1
  \item Option B: Get \$5, other gets \$5
\end{itemize}

Many chose B (fair). But when they could choose ``not to know'' what the other gets,
many chose ignorance and then A. They preferred not to know they were being unfair.

\paragraph{Implication.}
Utility includes self-image:
\begin{equation}
  U = U^{material} + U^{IDN}
\end{equation}

where $U^{IDN}$ is identity utility---the value of being (or seeing oneself as) a good person.

\paragraph{Temporal Complementarity.}
Identity creates intertemporal linkage:
\begin{equation}
  \frac{\partial^2 U}{\partial a_t \partial a_{t+1}} > 0
\end{equation}

Behaving as type $X$ today makes it easier to see oneself as type $X$ tomorrow.
This is the mechanism behind habit, character, and culture.

\subsubsection{Domain 3: Reference Dependence and Loss Aversion}

\paragraph{The Evidence.}
\citet{kahnemantversky1979} demonstrated that outcomes are evaluated
relative to a reference point, with losses weighted more than gains:
\begin{equation}
  v(x) = \begin{cases}
    x^\alpha & \text{if } x \geq 0 \\
    -\lambda (-x)^\beta & \text{if } x < 0
  \end{cases}
\end{equation}
with $\lambda \approx 2$ (loss aversion coefficient).

\paragraph{Example: The Endowment Effect.}
\citet{knetschetal1990} gave subjects mugs or chocolate bars.
Willingness to trade was far below 50\%---people valued what they had
more than equivalent items they didn't have.

Classical prediction: Indifferent subjects should trade 50\% of the time.

\paragraph{Implication.}
The reference point $r$ is part of context:
\begin{equation}
  U = U(x - r(\Psi))
\end{equation}

Reference points depend on:
\begin{itemize}
  \item Status quo (what I have)
  \item Expectations (what I thought I'd get) \citep{koszegirab2006}
  \item Social comparison (what others have)
  \item Aspirations (what I want to have)
\end{itemize}

All are context-dependent: $r = r(\Psi)$.

\subsubsection{Domain 4: Framing and Context Effects}

\paragraph{The Evidence.}
Identical choice problems elicit different responses depending on framing:
\begin{itemize}
  \item \textbf{Asian Disease Problem:} ``200 saved'' vs. ``400 die'' reverses preferences \citep{tverskykahneman1981}
  \item \textbf{Default effects:} Organ donation rates vary 10x based on opt-in vs. opt-out
  \item \textbf{Anchoring:} Arbitrary numbers influence willingness to pay
\end{itemize}

\paragraph{Example: Organ Donation Defaults.}
\citet{johnsonetal2003} showed:
\begin{itemize}
  \item Germany (opt-in): 12\% donation rate
  \item Austria (opt-out): 99\% donation rate
\end{itemize}

Same populations, similar cultures, radically different behavior---
entirely due to default setting.

\paragraph{Implication.}
Choices depend on how options are presented:
\begin{equation}
  \text{Choice}(A, B; \Psi_1) \neq \text{Choice}(A, B; \Psi_2)
\end{equation}

even when $A$ and $B$ are objectively identical.

Stable preferences would require: same options $\Rightarrow$ same choice.
This is empirically false.

\subsubsection{Domain 5: Preference Construction and Learning}

\paragraph{The Evidence.}
For many choices, preferences are not ``retrieved'' but ``constructed'' at decision time:
\begin{itemize}
  \item \textbf{Unfamiliar domains:} How much is a clean environment worth?
  \item \textbf{Complex tradeoffs:} How to weigh career vs. family?
  \item \textbf{New products:} Did you ``want'' a smartphone before it existed?
\end{itemize}

\paragraph{Example: Contingent Valuation.}
Willingness to pay for environmental goods varies wildly depending on:
\begin{itemize}
  \item Order of questions
  \item Starting point for bidding
  \item Whether WTP or WTA is elicited
  \item What comparison goods are mentioned
\end{itemize}

\paragraph{Implication.}
Preferences for complex goods are not stable attributes
but emerge from deliberation in context:
\begin{equation}
  U(x; \Psi) = f(\text{construction process}(\Psi))
\end{equation}

This is especially true for the ``higher'' FEPSDE dimensions:
emotional, social, identity, and ecological utility
are constructed through social processes, not simply discovered.

\subsubsection{Domain 6: Temporal Preferences and Present Bias}

\paragraph{The Evidence.}
A vast literature documents systematic deviations from exponential discounting
\citep{frederick2002time}:
\begin{itemize}
  \item \textbf{Present bias:} People overweight immediate rewards relative to 
        future rewards, even controlling for pure time preference
  \item \textbf{Preference reversals:} The same person prefers \$110 in 31 days 
        over \$100 in 30 days, but \$100 today over \$110 tomorrow
  \item \textbf{Commitment demand:} People pay to restrict their future choices \citep{ashraf2006,bryan2010}
  \item \textbf{Procrastination:} Costly tasks are perpetually postponed
\end{itemize}

\paragraph{The Time Inconsistency Problem.}
\citet{strotz1956} first identified that non-exponential discounting leads to
dynamically inconsistent preferences: optimal plans made today are not followed tomorrow.

\paragraph{The $\beta$-$\delta$ Model.}
\citet{laibson1997} formalized quasi-hyperbolic discounting:
\begin{equation}
  U_t = u(c_t) + \beta \sum_{\tau=1}^{T} \delta^\tau u(c_{t+\tau}) 
  \quad \text{where } \beta < 1
  \label{eq:beta-delta}
\end{equation}

The parameter $\beta$ captures present bias---the additional weight on immediate
consumption beyond standard exponential discounting $\delta$.

\paragraph{Empirical Estimates.}
Studies consistently find $\beta \approx 0.7$ and $\delta \approx 0.99$
\citep{frederick2002time,cohen2020measuring}. This implies:
\begin{itemize}
  \item Substantial present bias (30\% premium on immediate rewards)
  \item Modest long-run impatience
  \item Time-inconsistent behavior: plans made today are not followed tomorrow
\end{itemize}

\paragraph{Sophistication vs.\ Naïveté.}
\citet{odonoghue1999doing, odonoghue2001} distinguished:
\begin{itemize}
  \item \textbf{Sophisticates:} Know their $\beta < 1$, may seek commitment
  \item \textbf{Naïfs:} Believe $\hat{\beta} = 1$, perpetually surprised by own behavior
  \item \textbf{Partial naïfs:} $\beta < \hat{\beta} < 1$---underestimate but don't ignore bias
\end{itemize}

\paragraph{Framework Integration.}
Temporal preferences constitute the $\Psi_T$ dimension of context:
\begin{equation}
  \Psi_T = (\beta, \delta, \hat{\beta})
\end{equation}

Present bias implies that $C^*$ depends on the timing of decisions:
\begin{equation}
  C^*_{immediate} \neq C^*_{planned}
\end{equation}

This creates demand for commitment devices---institutional arrangements ($\Psi_I$)
that constrain future choices. Examples include:
\begin{itemize}
  \item Illiquid savings accounts (401k, Christmas clubs)
  \item Deadlines and penalties
  \item Default enrollment in beneficial programs \citep{madrian2001}
\end{itemize}

\citet{chetty2014} found that 85\% of savers are ``passive'' (follow defaults)
while only 15\% are ``active'' (respond to incentives)---demonstrating
the power of institutional design for present-biased agents.

\paragraph{Implications.}
Temporal context interacts with other $\Psi$ dimensions:
\begin{itemize}
  \item $\Psi_T \times \Psi_C$: Cognitive load worsens self-control
  \item $\Psi_T \times \Psi_I$: Institutions can provide commitment
  \item $\Psi_T \times \Psi_E$: Poverty may intensify present bias
\end{itemize}

\subsection{The Deeper Problem: Utility Interdependence}

\paragraph{Beyond Individual Violations.}
Each domain above could be treated as a ``behavioral anomaly''---
a deviation from rationality requiring a patch.

But the deeper issue is \emph{structural}:
utilities are interdependent in ways that classical theory cannot accommodate.

\paragraph{Four Types of Interdependence.}

\textbf{Cross-agent interdependence:}
\begin{equation}
  \frac{\partial U_i}{\partial x_j} \neq 0
\end{equation}
My utility depends on your outcome (social preferences).

\textbf{Temporal interdependence:}
\begin{equation}
  \frac{\partial U_t}{\partial a_{t-1}} \neq 0
\end{equation}
My current utility depends on my past actions (habits, identity).

\textbf{Intertemporal inconsistency:}
\begin{equation}
  U_t(c_t, c_{t+1}) \neq \delta \cdot U_{t+1}(c_{t+1}, c_{t+2})
\end{equation}
My time preferences are not constant (present bias).

\textbf{Context interdependence:}
\begin{equation}
  \frac{\partial U}{\partial \Psi} \neq 0
\end{equation}
My utility depends on the institutional environment (framing, norms).

\paragraph{The Complementarity Implication.}
All four types generate non-zero second derivatives:
\begin{equation}
  C_{ij} = \frac{\partial^2 U_i}{\partial x_i \partial x_j} \neq 0
\end{equation}

This is not a minor technical detail---it is the structural feature
that creates coordination problems, multiple equilibria, and path dependence.

\subsection{What Utility Maximization Gets Wrong}

\paragraph{Not Wrong, But Incomplete.}
The problem is not that utility maximization is false.
Given a utility function and constraints, maximization is a reasonable model of choice.

The problem is that the framework presupposes:
\begin{enumerate}
  \item The utility function is known and stable
  \item The constraints are exogenous
  \item Agents optimize independently
\end{enumerate}

When utilities are interdependent and context is endogenous,
``maximize utility'' doesn't define behavior---
it merely redescribes whatever happens as ``utility-maximizing.''

\paragraph{The Identification Problem.}
If utility can include anything (other-regarding preferences, identity, 
reference points, framing effects), then any behavior is ``rational.''
The theory becomes unfalsifiable.

\paragraph{The Equilibrium Selection Problem.}
With interdependent utilities, multiple equilibria exist.
Utility maximization says: at equilibrium, everyone is maximizing.
It doesn't say which equilibrium will emerge.

\paragraph{The Stability Problem.}
Utility maximization describes optimal response to given conditions.
It doesn't describe whether those conditions will persist---
whether the equilibrium is stable or fragile.

\subsection{Why These Deviations Cannot Be Learned from Text}

\paragraph{The Documentation Problem.}
The six domains above are well-documented in the behavioral economics literature.
Thousands of papers describe social preferences, loss aversion, framing effects,
and present bias. This documentation creates a puzzle: if LLMs are trained on
this literature, why can't they predict behavioral deviations?

\paragraph{The Peng et al. Finding.}
\citet{peng2025} found that LLM-based ``digital twins'' replicated only 50\% of
classic behavioral economics experiments. The pattern of failures is instructive:
\begin{itemize}
    \item \textbf{Attraction effect:} Humans show it; digital twins did not replicate
    \item \textbf{Compromise effect:} Weak in humans; not replicated by twins
    \item \textbf{Default effects:} Varied by context in humans; twins showed uniform patterns
\end{itemize}

\paragraph{The Core Problem: Propositional vs.\ Functional Knowledge.}
LLMs learn from text that ``loss aversion exists'' and ``$\lambda \approx 2$.''
This is \emph{propositional knowledge}---statements about facts.

But prediction requires \emph{functional knowledge}:
\begin{equation}
    \lambda(\Psi) = f(\Psi_C, \Psi_T, \Psi_E, \text{stake size}, \text{domain}, \ldots)
\end{equation}

When is $\lambda = 1.5$ versus $\lambda = 2.5$? When does loss aversion disappear entirely
\citep{list2003}? These questions cannot be answered from textual descriptions---
they require calibration on experimental variation.

\paragraph{Example: Social Preferences.}
The literature documents that people reject unfair ultimatum offers.
An LLM knows this. But can it predict:
\begin{itemize}
    \item Rejection rates in the Machiguenga (Peru) versus Pittsburgh?
    \item How rejection changes with stake size (\$10 vs.\ \$100 vs.\ \$10,000)?
    \item The effect of earned versus windfall endowments?
    \item Differences between student and representative samples?
\end{itemize}

These are exactly the questions that matter for prediction and policy.
The answers exist---scattered across hundreds of experiments.
But they are not in the text that LLMs read; they are in the data.

\paragraph{The Calibration Opportunity.}
This chapter's six domains are not merely ``anomalies'' to be catalogued.
They are \emph{calibration data} for $C^*(\Psi)$:

\begin{center}
\small
\begin{tabular}{lll}
\toprule
\textbf{Domain} & \textbf{What Varies} & \textbf{Calibration Target} \\
\midrule
Social preferences & Across cultures, stakes, merit & $C_{ij}(\Psi_S, \Psi_I)$ \\
Identity & Across groups, priming, salience & $U^{IDN}(\Psi_S, \text{group})$ \\
Reference dependence & Across domains, expectations & $r(\Psi)$, $\lambda(\Psi)$ \\
Framing & Across presentations, defaults & Frame sensitivity $= f(\Psi_C)$ \\
Preference construction & Across familiarity, complexity & Construction weight $= f(\Psi_K)$ \\
Temporal preferences & Across cultures, cognitive load & $\beta(\Psi_C, \Psi_T)$ \\
\bottomrule
\end{tabular}
\end{center}

Each experimental finding that documents context-dependence adds information
to the functional form of $C^*(\Psi)$. The variation that makes behavioral
economics messy is precisely what makes calibration possible.

\paragraph{From Anomalies to Architecture.}
The behavioral economics literature spent decades documenting deviations from
classical rationality. This was necessary and valuable---establishing \emph{that}
people deviate. But the next step is to predict \emph{when} and \emph{how much}.

Section~4 develops the empirical foundations for this prediction.
Section~4.X explains why calibration on experimental data succeeds where
LLM-based simulation fails. The six domains above are not problems to be
solved but resources to be exploited.

\subsection{The Required Extension}

\paragraph{From Optimization to Coherence.}
The limits of utility maximization suggest a different criterion:
not ``is each agent optimizing?'' but ``is the system of utilities coherent?''

Coherence means:
\begin{itemize}
  \item Expectations match outcomes
  \item Institutions support behavior
  \item Norms align with incentives
  \item Identity is consistent with action
\end{itemize}

\paragraph{The Framework Extension.}
\begin{enumerate}
  \item Allow $C_{ij} \neq 0$: Interdependent utilities
  \item Allow $\Psi_{t+1} \neq \Psi_t$: Endogenous context
  \item Define $K$: Coherence as structural stability
  \item Define $Q$: Quality as welfare at equilibrium
  \item Analyze stability: When does $K \approx 1$ persist?
\end{enumerate}

This is what the following sections develop.

\paragraph{The Methodological Shift.}
The extension requires not just theoretical generalization but a methodological shift:
from \emph{deriving} behavior from axioms to \emph{calibrating} behavior on data.
The six domains above generate predictions only when their parameters are estimated
from experimental variation.

%------------------------------------------------------------------------------
% SUMMARY
%------------------------------------------------------------------------------
\section{Summary}
\label{sec:ch3-summary}

\begin{tcolorbox}[
    colback=blue!5!white,
    colframe=blue!75!black,
    title={\textbf{Chapter 3 Summary: Limits of Classical Utility}}
]

\textbf{Core Insight:} Classical utility maximization under stable preferences excludes six fundamental behavioral domains that require interdependence ($C_{ij} \neq 0$) and endogenous context ($\Psi_{t+1} \neq \Psi_t$).

\textbf{Key Results:}
\begin{enumerate}[nosep]
    \item \textbf{Reference Dependence:} Utility depends on reference points, not just final states ($\lambda \approx 2.25$)
    \item \textbf{Social Preferences:} Fairness considerations violate self-interest ($\alpha_{fairness} > 0$)
    \item \textbf{Time Inconsistency:} Present bias creates preference reversals ($\beta < 1$)
    \item \textbf{Limited Attention:} Bounded rationality constrains optimization
    \item \textbf{Identity Effects:} Self-concept influences choices
    \item \textbf{Context Dependence:} Framing and salience matter
\end{enumerate}

\textbf{Transition:} Chapter 4 develops the empirical foundations for measuring these effects.
\end{tcolorbox}

\subsection{What Comes Next}
\label{sec:ch3-next}

With the limitations of classical utility established, subsequent chapters build the empirical foundation:

\begin{enumerate}[nosep]
    \item \textbf{Chapter 4:} Develops empirical foundations---from theory to measurement
    \item \textbf{Chapter 4x:} Explains why calibration succeeds where LLM-based simulation fails
    \item \textbf{Chapter 5:} Introduces complementarity as the structural principle for interdependence
\end{enumerate}

%%%%%%%%%%%%%%%%%%%%%%%%%%%%%%%%%%%%%%%%%%%%%%%%%%%%%%%%%%%%%%%%%%%%%%%%%%%%%%%
% READING PATH BOX
%%%%%%%%%%%%%%%%%%%%%%%%%%%%%%%%%%%%%%%%%%%%%%%%%%%%%%%%%%%%%%%%%%%%%%%%%%%%%%%
\begin{tcolorbox}[
    colback=green!5!white,
    colframe=green!75!black,
    title={\textbf{Reading Path}},
    fonttitle=\bfseries
]
\begin{tabular}{@{}ll@{}}
\textbf{Previous:} & Chapter 2 \\
\textbf{Next:} & Chapter 4 \\
\textbf{Deep Dive:} & See Appendix G (Glossary) for definitions \\
\end{tabular}

\medskip
\textit{For readers short on time: Focus on the Intuition Box and Summary section.}
\end{tcolorbox}

